\section{Proyecciones y denominadores}
\label{sec:proyecciones}

\fuente{B, §5.1, especialmente Lema~5.1, p.~53, y §§5.2--5.3,
pp.~54--57; A, Lema~6.3, p.~22, y §8, pp.~35--42. El ejemplo de
dimensión dos y el lema local se demuestran aquí.}

\subsection{Un operador integral con un proyector no integral}

Sea $p$ primo, sea $V=\Q^2$ y sea $M=\Z^2$ su retículo habitual.
Consideremos
\begin{equation}
 T=\begin{pmatrix}0&1\\0&p\end{pmatrix},
 \qquad T(x,y)=(y,py).
 \label{eq:operador-t}
\end{equation}
La matriz tiene entradas enteras, así que $T(M)\subseteq M$. Sus dos
autovalores son $0$ y $p$, con autoespacios
\[
 V_0=\Q(1,0),\qquad V_p=\Q(1,p).
\]
Todo vector se descompone de manera única como
\begin{equation}
 (x,y)=\left(x-\frac yp\right)(1,0)+\frac yp(1,p).
 \label{eq:descomposicion}
\end{equation}
Por tanto, el proyector sobre $V_p$ a lo largo de $V_0$ es
\begin{equation}
 P=\frac{T}{p}
   =\begin{pmatrix}0&1/p\\0&1\end{pmatrix}.
 \label{eq:proyector-p}
\end{equation}
La igualdad $T^2=pT$ comprueba que $P^2=P$. Sin embargo,
$P(0,1)=(1/p,1)\notin M$. El operador $T$ es integral y diagonalizable
sobre $\Q$, pero su proyector espectral no preserva $M$.

La diferencia entre los autovalores es $p$, que no es una unidad sobre
$\Z_p$. La fórmula $P=(T-0I)/(p-0)$ muestra exactamente dónde aparece
el denominador. Módulo $p$, los autovalores coinciden y
\[
 \overline T=\begin{pmatrix}0&1\\0&0\end{pmatrix}
\]
es un operador nilpotente no nulo. La separación de autoespacios sobre
el cuerpo no se conserva como descomposición del retículo.

\subsection{Dos cocientes que conviene distinguir}

Los retículos que están dentro de cada autoespacio son
\[
 M\cap V_0=\Z(1,0),\qquad M\cap V_p=\Z(1,p).
\]
Su suma es
\[
 M'=(M\cap V_0)\oplus(M\cap V_p)
    =\{(a+b,pb):a,b\in\Z\}=\Z\times p\Z.
\]
Así, $[M:M']=p$. Este índice mide el coste de sustituir $M$ por un
retículo que sí se descompone en los autoespacios.

El proyector, en cambio, produce
\[
 PM=\frac1p\Z(1,p),
 \qquad
 \frac{PM}{M\cap V_p}
 =\frac{(1/p)\Z(1,p)}{\Z(1,p)}\simeq\Z/p\Z.
\]
Este segundo cociente mide cuánto amplía el proyector el retículo de
$V_p$. Es el objeto que aparece directamente en el Lema~5.1 de B.
Ambos índices valen $p$ en el ejemplo, pero sus definiciones son distintas.

\subsection{El lema local de denominadores}

\begin{lema}
Sea $\mathcal O$ un anillo de valoración discreta con cuerpo de fracciones
$K$ y valoración $v$ normalizada por $v(\varpi)=1$. Sea $M$ un retículo
pleno en $V=V_1\oplus V_2$, es decir, un $\mathcal O$-submódulo
finitamente generado que genera $V$ sobre $K$, y sea $e$ el proyector sobre $V_1$.
Supongamos que un operador $T$ preserva $M$, actúa por un escalar
$a\in\mathcal O\setminus\{0\}$ sobre $V_1$ y por cero sobre $V_2$.
Entonces
\begin{equation}
 aeM\subseteq M\cap V_1,
 \qquad
 \length_{\mathcal O}\frac{eM}{M\cap V_1}
 \le (\dim_K V_1)v(a).
 \label{eq:lema-denominadores}
\end{equation}
\end{lema}

\begin{proof}
Se tiene $T=ae$ en $V$. Para $m\in M$, el vector $aem=Tm$ pertenece
a $M$ y a $V_1$. Además, $M\cap V_1\subseteq eM$, pues $e$ es la
identidad en $V_1$. El cociente de \eqref{eq:lema-denominadores} es,
por tanto, un cociente de $eM/aeM$.

El módulo $eM$ es finitamente generado, libre de torsión y genera $V_1$;
sobre un anillo de valoración discreta es libre de rango $d_1=\dim_K V_1$.
Por ello
\[
 \length_{\mathcal O}(eM/aeM)
 =d_1\length_{\mathcal O}(\mathcal O/a\mathcal O)=d_1v(a).
\]
La longitud no aumenta al pasar a un cociente, lo que prueba la cota.
\end{proof}

Para \eqref{eq:operador-t}, después de extender de $\Z$ a $\Z_p$, se
tiene $a=p$, $d_1=1$ y longitud $1$. La desigualdad es una igualdad.
En una valoración discreta, \emph{longitud} cuenta factores isomorfos al
cuerpo residual; sobre $\Z_p$, un módulo finito de longitud $r$ tiene
orden $p^r$.

\subsection{Del ejemplo a los periodos automorfos}

Los espacios automorfos utilizados por las fuentes admiten descomposiciones
en sistemas de autovalores de Hecke sobre un cuerpo de coeficientes.
Los operadores de Hecke preservan estructuras integrales, pero proyectar
en los autoespacios elegidos puede introducir denominadores. Si dos
sistemas de autovalores son congruentes módulo un primo de coeficientes,
un operador que anula uno puede actuar por un escalar divisible por ese
primo sobre el otro. El lema cuantifica el coste de dividir por tal escalar.

La conexión con la comparación de pendientes puede resumirse sin ocultar
sus objetos. Para un cuerpo de coeficientes $K$, las fuentes construyen
un emparejamiento normalizado
\[
 P_\alpha:U_\alpha^\circ\times W\longrightarrow K
\]
y retículos $M'_\alpha\subset U_\alpha^\circ$, $M_W\subset W$.
La divisibilidad da un ideal fraccionario $\mathfrak a_\alpha$ con
\[
 P_\alpha(M'_\alpha,M_W)\subseteq\mathfrak a_\alpha.
\]
Al definir
\[
 N_\alpha=\mathfrak a_\alpha^{-1}M'_\alpha,
 \qquad C=\Hom_{\OK}(M_W,\OK),
\]
se obtiene una aplicación integral
\[
 b_\alpha:N_\alpha\longrightarrow C,
 \qquad b_\alpha(f)(g)=P_\alpha(f,g).
\]
Ampliar el retículo fuente consume la divisibilidad disponible; estimar
las normas de $b_\alpha$ permite comparar esa ganancia con el coste
complejo. Los denominadores de las proyecciones restan parte de la ganancia
y deben entrar en la estimación. Véanse B, pp.~29 y 49--56, y A,
Lema~6.3 y §8.

\subsection{Por qué el siguiente paso es cohomológico}

En la construcción profunda de las fuentes, las congruencias entre formas
automorfas se estudian mediante extensiones de sus representaciones de
Galois. Las condiciones locales de esas extensiones llevan, entre otros
objetos, a los grupos
\begin{equation}
 S_\ell(E,j)=\ker\left(
 H^1(G_{E,S_\ell},\Q_\ell/\Z_\ell(-j))
 \longrightarrow\prod_{w\mid\ell}
 H^1(E_w,\Q_\ell/\Z_\ell(-j))\right),
 \label{eq:grupo-selmer}
\end{equation}
donde $j\ge1$ y $S_\ell$ es el conjunto de lugares de $E$ sobre $\ell$.
El signo del twist es parte de la definición. Si el grupo es finito, su
longitud multiplicada por $\log\ell$ es el logaritmo de su orden.

Para el argumento no basta saber que un grupo aislado es finito. Se
necesita controlar sumas de estas longitudes sobre primos y twists, con
uniformidad al variar los campos auxiliares. A, §8, y B, §5, formulan
esas cotas y la selección de una ventana común. Esta nota explica por
qué hacen falta; su construcción sigue siendo un resultado de las fuentes
pendiente de desarrollar en el cuaderno.

\begin{ejercicio}
Sustituya $p$ por un entero $d\ge2$ en \eqref{eq:operador-t}.
Calcule $[M:M']$ y $PM/(M\cap V_d)$. Después localice en un primo
$\ell$: ¿qué longitud obtiene y para qué primos el proyector ya es integral?
\emph{Comprobación:} el índice es $d$, el cociente es $\Z/d\Z$ y la
longitud local es $v_\ell(d)$.
\end{ejercicio}
